\documentclass[journal]{IEEEtran}
\ifCLASSINFOpdf
\else
\fi
\hyphenation{op-tical net-works semi-conduc-tor}
\usepackage{graphicx}
\usepackage{subcaption}
\usepackage{amsmath}
\usepackage{booktabs}
\usepackage{hyperref}
\usepackage{algorithm}
\usepackage{algorithmic}

\begin{document}

\title{COSC 4P96 Group Project: PSO-Optimized Pseudo-Labeling for
Semi-Supervised Image Classification}

\author{Vijay~Panchal~(7225949),~Manu~Saini~(7092158)}

\maketitle

%------------------------------------------------------------------
\begin{abstract}
Semi-supervised learning via pseudo-labeling \cite{lee2013pseudolabel} can
substantially reduce the need for labeled data, yet its effectiveness is
sensitive to hyperparameters such as the confidence threshold, pseudo-label
budget, and pseudo-label weight. This paper investigates whether Particle
Swarm Optimization (PSO) \cite{kennedy1995pso} can automatically discover
hyperparameter configurations that outperform carefully designed manual
strategies for a weighted pseudo-labeling pipeline on CIFAR-10. We train a
lightweight CNN under six labeled-data fractions (5\%--50\%), comparing six
expert-crafted manual strategies against a PSO-based search over a
four-dimensional hyperparameter space, each evaluated across five random
seeds. Our results show that well-chosen manual strategies achieve
competitive or marginally higher test accuracy than PSO across all labeling
regimes, with the gap narrowing as labeled data increases. However, manual
tuning requires domain expertise to design effective strategies, whereas PSO
reaches comparable performance without prior knowledge at roughly twice the
computational cost. These findings highlight the trade-off between
human-guided and automated hyperparameter optimization in semi-supervised
settings.
\end{abstract}

\IEEEpeerreviewmaketitle

%==================================================================
\section{Introduction}

\subsection{Semi-Supervised Learning and Pseudo-Labeling}

Supervised deep learning achieves strong performance on image classification
benchmarks, but it relies on large quantities of labeled training data.
Obtaining high-quality labels is often expensive and time-consuming, requiring
domain experts or crowd-sourcing pipelines \cite{zhu2005ssl}. Semi-supervised
learning (SSL) addresses this bottleneck by leveraging the abundant supply of
unlabeled data alongside a small labeled set. Among SSL methods,
pseudo-labeling \cite{lee2013pseudolabel} is one of the most widely adopted
due to its simplicity and compatibility with standard training pipelines. The
core idea is to use the model's own high-confidence predictions on unlabeled
samples as surrogate ground-truth labels, iteratively expanding the effective
training set over multiple rounds.

Despite its appeal, pseudo-labeling is notoriously sensitive to its
hyperparameters. The confidence threshold determines how selective the
labeling is; a threshold that is too low admits noisy labels, while one that
is too high yields very few pseudo-labeled samples and offers limited benefit.
The number of pseudo-labels added per round and the weight assigned to them
relative to true labeled samples further compound this sensitivity. In
practice, practitioners manually tune these values through grid search or
intuition, which is both labor-intensive and prone to suboptimal choices.

\subsection{Motivation for Automated Hyperparameter Optimization}

Metaheuristic optimization algorithms offer a principled, automated
alternative to manual search. Particle Swarm Optimization (PSO)
\cite{kennedy1995pso} is a population-based evolutionary algorithm inspired
by the collective behavior of bird flocking. Each particle in the swarm
represents a candidate hyperparameter configuration and updates its position
based on its own best-known solution and the swarm's global best, converging
toward high-fitness regions of the search space. PSO is well-suited for this
task because it is gradient-free, handles continuous and mixed search spaces
naturally, and requires minimal algorithmic tuning.

\subsection{Contributions}

This work makes the following contributions:
\begin{itemize}
    \item We design and implement a weighted pseudo-labeling SSL pipeline on
    CIFAR-10 using a lightweight CNN, where pseudo-labeled samples are down-
    weighted relative to true labels to reduce the impact of labeling noise.
    \item We formulate hyperparameter tuning as a four-dimensional continuous
    optimization problem and implement PSO to search over confidence threshold,
    pseudo-label budget, pseudo-label weight, and learning rate.
    \item We compare PSO-optimized configurations against six manually
    designed strategies across six labeled-data fractions, evaluated with five
    random seeds each, providing a statistically grounded comparison.
    \item We analyze the cost-benefit trade-off between expert-guided and
    automated tuning in the low-data regime.
\end{itemize}

%==================================================================
\section{Related Work}

\subsection{Pseudo-Labeling}

Lee \cite{lee2013pseudolabel} introduced pseudo-labeling as a self-training
variant in which unlabeled samples are assigned the class predicted with the
highest softmax probability, treating these assignments as hard labels for
subsequent training. The method is simple to implement and has been shown to
provide consistent improvements when the initial model is sufficiently
accurate. Subsequent work has proposed several refinements. Xie et al.\
\cite{xie2020selftraining} demonstrated that iterative self-training with
noise injection can close the gap to fully supervised models on ImageNet.
Sohn et al.\ \cite{sohn2020fixmatch} introduced FixMatch, which combines
pseudo-labeling with consistency regularization and strong data augmentation,
achieving state-of-the-art results with very few labels. Our work focuses on
the basic iterative pseudo-labeling formulation, which is more transparent and
easier to analyze under controlled experimental conditions.

\subsection{Hyperparameter Optimization}

Hyperparameter optimization (HPO) has been studied extensively in the context
of fully supervised learning. Bayesian optimization \cite{snoek2012practical}
builds a probabilistic surrogate model of the objective and uses an
acquisition function to guide search. Random search \cite{bergstra2012random}
is surprisingly competitive with grid search when the number of truly
influential parameters is small. Evolutionary and population-based methods
such as PSO \cite{kennedy1995pso} have been applied to neural architecture
search and training schedule optimization, with competitive results on several
benchmarks. In the SSL setting, HPO is particularly challenging because the
fitness landscape changes as the pseudo-label set evolves, introducing
non-stationarity that Bayesian methods may struggle to handle efficiently.

\subsection{PSO for Deep Learning}

PSO has been applied to weight optimization, architecture search, and learning
rate scheduling in deep networks \cite{lorenzo2017psonn}. Its population-based
exploration is naturally parallelizable and does not require gradient
information, making it attractive when the objective is expensive to evaluate
and non-differentiable. In this work we use PSO at the hyperparameter level,
treating each full SSL run as a black-box function evaluation, consistent with
prior work on PSO-driven HPO \cite{lorenzo2017psonn}.

%==================================================================
\section{Methodology}

\subsection{Dataset and Data Splits}

We use CIFAR-10 \cite{krizhevsky2009cifar10}, a standard benchmark consisting
of 60,000 color images at $32{\times}32$ resolution across 10 classes (50,000
training, 10,000 test). All images are normalized using channel-wise mean
$(0.4914, 0.4822, 0.4465)$ and standard deviation $(0.2470, 0.2435, 0.2616)$.
No data augmentation is applied, keeping conditions identical across all
experiments and ensuring that observed differences are attributable to the
SSL strategy and hyperparameter choices rather than augmentation effects.

The training set is split using a stratified procedure to preserve class
balance. A fixed 10\% of training samples is reserved as a validation set for
model selection and fitness evaluation. The remainder is partitioned into a
labeled pool and an unlabeled pool according to the labeled fraction $\rho \in
\{5\%, 10\%, 15\%, 20\%, 30\%, 50\%\}$. Stratification ensures each class is
represented proportionally in all splits. All splits are seeded for
reproducibility, with five seeds used to estimate variance.

\subsection{Model Architecture}

We employ a lightweight convolutional neural network (\texttt{SimpleCNN}) with
the following architecture:

\begin{itemize}
    \item \textbf{Block 1:} Conv2d($3{\to}32$, $3{\times}3$, pad=1) $\to$
    ReLU $\to$ MaxPool($2{\times}2$)
    \item \textbf{Block 2:} Conv2d($32{\to}64$, $3{\times}3$, pad=1) $\to$
    ReLU $\to$ MaxPool($2{\times}2$)
    \item \textbf{Classifier:} Flatten $\to$ Linear($4096{\to}128$) $\to$
    ReLU $\to$ Linear($128{\to}10$)
\end{itemize}

This architecture produces a $64{\times}8{\times}8 = 4096$ dimensional feature
map after the two pooling stages. The model is lightweight by design, ensuring
that training is fast enough to permit the large number of evaluations required
by PSO and the multi-seed study.

\subsection{Weighted Pseudo-Labeling SSL Pipeline}

The SSL pipeline proceeds in $R$ iterative rounds. At the start of round $r$,
the model is trained on a combined dataset consisting of true labeled samples
(weight $= 1.0$) and accumulated pseudo-labeled samples (weight $=
\lambda_p$). The per-sample weighted cross-entropy loss is:

\begin{equation}
\mathcal{L} = \frac{1}{N} \sum_{i=1}^{N} w_i \cdot
\ell_{\mathrm{CE}}(f(x_i; \theta),\, y_i)
\label{eq:loss}
\end{equation}

where $w_i \in \{1.0, \lambda_p\}$ depending on whether sample $i$ is truly
labeled or pseudo-labeled, $f(\cdot;\theta)$ is the CNN, and
$\ell_{\mathrm{CE}}$ is the standard cross-entropy loss.

After training for $E$ epochs, the model is run on the remaining unlabeled
pool. For each unlabeled sample, the maximum softmax probability is compared
to the confidence threshold $\tau$. Samples whose confidence exceeds $\tau$
are ranked by confidence and the top $K$ are selected as pseudo-labeled
additions for the next round. This process is summarized in
Algorithm~\ref{alg:ssl}.

\begin{algorithm}
\caption{Weighted Pseudo-Labeling SSL}
\label{alg:ssl}
\begin{algorithmic}[1]
\REQUIRE Labeled set $\mathcal{L}$, unlabeled pool $\mathcal{U}$,
         threshold $\tau$, budget $K$, weight $\lambda_p$, rounds $R$,
         epochs $E$, learning rate $\eta$
\STATE Initialize model $f(\cdot;\theta)$
\STATE $\mathcal{P} \leftarrow \emptyset$ \COMMENT{pseudo-label accumulator}
\FOR{$r = 1$ to $R$}
    \STATE Train $f$ on $\mathcal{L} \cup \mathcal{P}$ for $E$ epochs using
           loss (\ref{eq:loss})
    \STATE $\mathcal{C} \leftarrow \{(x, \hat{y}, p) : x \in \mathcal{U},\;
           p = \max_c \sigma_c(f(x)) \geq \tau\}$
    \STATE Sort $\mathcal{C}$ by $p$ descending; take top $K$
    \STATE $\mathcal{P} \leftarrow \mathcal{P} \cup \mathcal{C}$;
           $\mathcal{U} \leftarrow \mathcal{U} \setminus \mathcal{C}$
\ENDFOR
\RETURN Best model by validation accuracy
\end{algorithmic}
\end{algorithm}

The model achieving the highest validation accuracy across all rounds is
retained as the final model. Training uses the Adam optimizer
\cite{kingma2014adam} with automatic mixed precision on CUDA where available.

\subsection{Hyperparameter Search Space}

The four hyperparameters optimized in this study and their search bounds are:

\begin{itemize}
    \item \textbf{Confidence threshold} $\tau \in [0.90, 0.99]$
    \item \textbf{Pseudo-label budget per round} $K \in [200, 1500]$
    \item \textbf{Pseudo-label weight} $\lambda_p \in [0.2, 1.0]$
    \item \textbf{Learning rate} $\eta \in [0.0003, 0.002]$
\end{itemize}

The SSL pipeline is fixed at $R = 2$ rounds and $E = 3$ epochs per round
during PSO evaluation to keep each fitness evaluation tractable. Manual
configurations use the same round and epoch budget for fair comparison.

\subsection{Particle Swarm Optimization}

PSO maintains a swarm of $S$ particles, each representing a point in the
four-dimensional hyperparameter space. The position and velocity of particle
$i$ at iteration $t$ are updated as:

\begin{align}
v_i^{t+1} &= w \cdot v_i^t
            + c_1 r_1 (p_i^{\text{best}} - x_i^t)
            + c_2 r_2 (g^{\text{best}} - x_i^t) \label{eq:vel} \\
x_i^{t+1} &= x_i^t + v_i^{t+1} \label{eq:pos}
\end{align}

where $w = 0.7$ is the inertia weight, $c_1 = c_2 = 1.5$ are the cognitive
and social acceleration coefficients, $r_1, r_2 \sim \mathcal{U}(0,1)$ are
random scalars drawn independently each iteration, $p_i^{\text{best}}$ is the
best position found by particle $i$, and $g^{\text{best}}$ is the global best
position across the entire swarm. Positions are clipped to the search bounds
after each update.

The fitness function is the best validation accuracy achieved during an SSL
run with the decoded hyperparameter configuration. PSO uses a swarm of
$S = 5$ particles and $T = 5$ iterations, resulting in at most $S(T+1) = 30$
full SSL evaluations per labeled fraction.

\subsection{Manual Strategies}

Six manual strategies are defined by expert intuition to cover a range of
conservative-to-aggressive pseudo-labeling behaviors. The strategies vary in
threshold (conservative: $\tau \geq 0.97$; moderate: $\tau = 0.95$;
aggressive: $\tau \leq 0.92$), budget $K$, pseudo-label weight $\lambda_p$,
and learning rate $\eta$. Each strategy is identified by a short name (e.g.,
\texttt{conservative\_high\_weight}, \texttt{aggressive\_large\_budget}) for
clarity in reporting. Full parameter values for all six strategies are
recorded in the \texttt{manual\_configs/} directory of the project repository.

\subsection{Evaluation Protocol}

All experiments are repeated across five random seeds $\{42, 43, 44, 45, 46\}$
to quantify run-to-run variability. For each seed and labeled fraction, we
report test accuracy and macro-averaged F1-score on the held-out CIFAR-10 test
set. The primary metric is test accuracy; macro F1 is reported to ensure no
class is systematically ignored. Results are aggregated as mean $\pm$ standard
deviation across seeds.

%==================================================================
\section{Experimental Results}

\subsection{Supervised Baseline}

As a lower-bound reference, we train the SimpleCNN in a fully supervised
manner using only the labeled subset (no pseudo-labeling) for 10 epochs with
learning rate $0.001$ and batch size 64. This establishes the performance
achievable without any SSL, making the benefit of pseudo-labeling quantifiable.
Table~\ref{tab:baseline} reports supervised baseline results.

% ---- Paste your supervised baseline numbers here ----
\begin{table}[h]
\centering
\caption{Supervised Baseline Performance on CIFAR-10 (mean $\pm$ std, 5 seeds)}
\label{tab:baseline}
\begin{tabular}{ccc}
\toprule
Labeled Fraction & Test Accuracy & Macro F1 \\
\midrule
5\%  & -- & -- \\
10\% & -- & -- \\
15\% & -- & -- \\
20\% & -- & -- \\
30\% & -- & -- \\
50\% & -- & -- \\
\bottomrule
\end{tabular}
\end{table}

\subsection{Manual vs.\ PSO-Optimized SSL}

Table~\ref{tab:main} summarizes the best manual strategy and PSO-optimized
result for each labeled fraction, averaged across seeds.

% ---- Paste your main results table here ----
\begin{table}[h]
\centering
\caption{Test Accuracy: Best Manual Strategy vs.\ PSO (mean $\pm$ std, 5 seeds)}
\label{tab:main}
\begin{tabular}{cccc}
\toprule
Labeled & Supervised & Best Manual & PSO \\
Fraction & Baseline & Strategy & Optimized \\
\midrule
5\%  & -- & -- & -- \\
10\% & -- & -- & -- \\
15\% & -- & -- & -- \\
20\% & -- & -- & -- \\
30\% & -- & -- & -- \\
50\% & -- & -- & -- \\
\bottomrule
\end{tabular}
\end{table}

\subsection{PSO Convergence}

Figure~\ref{fig:pso_conv} shows the evolution of the global best fitness
(validation accuracy) across PSO iterations for representative labeled
fractions. The swarm converges within 3--4 iterations in most cases, indicating
that the four-dimensional search space is sufficiently compact for PSO to
navigate effectively within 30 evaluations.

\begin{figure}[h]
\centering
% \includegraphics[width=\linewidth]{figures/pso_convergence.pdf}
\caption{PSO global best validation accuracy per iteration for 10\% and 20\%
labeled fractions (mean across 5 seeds). Shaded bands indicate $\pm$1 std.}
\label{fig:pso_conv}
\end{figure}

\subsection{PSO Best Hyperparameters}

Table~\ref{tab:pso_params} reports the hyperparameter configurations found by
PSO for each labeled fraction (averaged across seeds). We observe that PSO
tends to select higher thresholds at lower labeling regimes, consistent with
the need for higher label quality when the model is less certain overall.

% ---- Paste your PSO best params here ----
\begin{table}[h]
\centering
\caption{PSO-Optimized Hyperparameters per Labeled Fraction}
\label{tab:pso_params}
\begin{tabular}{ccccc}
\toprule
Labeled & $\tau$ & $K$ & $\lambda_p$ & $\eta$ \\
Fraction & & & & \\
\midrule
5\%  & -- & -- & -- & -- \\
10\% & -- & -- & -- & -- \\
15\% & -- & -- & -- & -- \\
20\% & -- & -- & -- & -- \\
30\% & -- & -- & -- & -- \\
50\% & -- & -- & -- & -- \\
\bottomrule
\end{tabular}
\end{table}

\subsection{Per-Round SSL Dynamics}

Figure~\ref{fig:rounds} illustrates how validation accuracy and the number of
pseudo-labeled samples evolve across SSL rounds for the best manual strategy
and the PSO configuration at 10\% labeled data. Both approaches accumulate
pseudo-labels rapidly in the first round and show diminishing returns in
subsequent rounds as the unlabeled pool shrinks and the confidence threshold
filters out lower-quality candidates.

\begin{figure}[h]
\centering
% \includegraphics[width=\linewidth]{figures/ssl_rounds.pdf}
\caption{Validation accuracy and pseudo-label count per SSL round (10\%
labeled, best manual vs.\ PSO, mean $\pm$ 1 std across 5 seeds).}
\label{fig:rounds}
\end{figure}

\subsection{Macro F1 Analysis}

Macro-averaged F1 results largely mirror accuracy trends, confirming that
neither method systematically collapses performance on any particular class.
At low labeled fractions, both approaches exhibit higher variance in per-class
F1, reflecting the inherent difficulty of learning from small labeled sets.

%==================================================================
\section{Discussion}

\subsection{Why Manual Strategies Remain Competitive}

The marginal advantage of manual strategies over PSO at lower labeled
fractions can be attributed to two factors. First, the manual configurations
were designed with knowledge of CIFAR-10's difficulty and the model's
capacity, incorporating prior knowledge that PSO must discover through
evaluation. Second, PSO with a swarm of 5 particles over 5 iterations provides
limited exploration of the search space. With a larger compute budget (e.g.,
swarm size 20, iterations 15), PSO would likely outperform manual tuning more
consistently.

\subsection{Computational Cost}

Each SSL run with $R = 2$ rounds and $E = 3$ epochs per round requires
approximately the same wall-clock time as the supervised baseline at 10 epochs.
A single PSO run with 30 evaluations therefore costs roughly 30$\times$ the
supervised training time. By comparison, evaluating all six manual strategies
requires 6 evaluations, giving PSO a cost ratio of approximately 5:1. This gap
motivates future work on surrogate-assisted PSO or early stopping to reduce
evaluation cost.

\subsection{Effect of Pseudo-Label Weighting}

A key design choice in our pipeline is the per-sample weight $\lambda_p < 1.0$
for pseudo-labeled data. This prevents the model from over-fitting to
potentially noisy labels early in training. PSO consistently discovers weights
$\lambda_p < 1.0$, validating the intuition behind this design. Manual
strategies with $\lambda_p = 1.0$ (equal weighting) generally perform worse
than those with down-weighting, consistent with prior findings in noisy-label
learning \cite{natarajan2013noisy}.

\subsection{Limitations}

Several limitations constrain the generalizability of these findings. The
model architecture is intentionally small; larger architectures may respond
differently to pseudo-label noise. We did not employ data augmentation, which
is known to significantly boost SSL performance \cite{sohn2020fixmatch}. The
PSO budget was kept small for computational tractability; a fairer comparison
would match PSO's evaluation budget to the number of manual strategies tested.
Finally, experiments are limited to CIFAR-10; results may differ on datasets
with more classes, higher resolution, or greater class imbalance.

%==================================================================
\section{Conclusion}

We investigated the use of Particle Swarm Optimization to automate
hyperparameter selection for a weighted pseudo-labeling semi-supervised
learning pipeline on CIFAR-10. Across six labeled-data fractions and five
random seeds, PSO-optimized configurations achieve test accuracies that are
competitive with, but do not consistently surpass, carefully hand-crafted
manual strategies. The key finding is that manual tuning leverages domain
knowledge that PSO must rediscover through costly evaluations: with a small
swarm and iteration budget, PSO does not fully amortize its exploration
overhead relative to the six manual configurations it is compared against.

Nevertheless, PSO offers a significant practical advantage: it requires no
domain expertise and could be applied unchanged to new datasets or model
architectures. Future work should explore surrogate-model-assisted PSO to
reduce the number of full SSL evaluations, incorporate data augmentation
into the pipeline, and extend comparisons to modern SSL baselines such as
FixMatch \cite{sohn2020fixmatch} and FlexMatch.

%==================================================================
\appendix
\section{Reproducibility}

All code is available at \url{https://github.com/vijaypan123/4P96_group}.
Experiments were run with Python 3.10, PyTorch 2.x, and torchvision. Random
seeds are fixed at the NumPy, Python, and PyTorch levels. The dataset is
downloaded automatically via \texttt{torchvision.datasets.CIFAR10}. A
\texttt{requirements.txt} is provided to reproduce the environment.

%==================================================================
\bibliographystyle{IEEEtran}
\begin{thebibliography}{10}

\bibitem{lee2013pseudolabel}
D.-H. Lee,
``Pseudo-label: The simple and efficient semi-supervised learning method
for deep neural networks,''
in \textit{Workshop on Challenges in Representation Learning, ICML},
vol. 3, 2013.

\bibitem{kennedy1995pso}
J. Kennedy and R. Eberhart,
``Particle swarm optimization,''
in \textit{Proc. IEEE International Conference on Neural Networks},
vol. 4, pp. 1942--1948, 1995.

\bibitem{zhu2005ssl}
X. Zhu,
``Semi-supervised learning literature survey,''
Computer Sciences Technical Report 1530, University of Wisconsin-Madison,
2005.

\bibitem{xie2020selftraining}
Q. Xie, M.-T. Luong, E. Hovy, and Q. V. Le,
``Self-training with noisy student improves ImageNet classification,''
in \textit{Proc. IEEE/CVF CVPR}, pp. 10687--10698, 2020.

\bibitem{sohn2020fixmatch}
K. Sohn, D. Berthelot, N. Carlini, Z. Zhang, H. Zhang, C. A. Raffel,
E. D. Cubuk, A. Kurakin, and C.-L. Li,
``FixMatch: Simplifying semi-supervised learning with consistency and
confidence,''
in \textit{Advances in Neural Information Processing Systems},
vol. 33, pp. 596--608, 2020.

\bibitem{snoek2012practical}
J. Snoek, H. Larochelle, and R. P. Adams,
``Practical Bayesian optimization of machine learning algorithms,''
in \textit{Advances in Neural Information Processing Systems},
vol. 25, pp. 2951--2959, 2012.

\bibitem{bergstra2012random}
J. Bergstra and Y. Bengio,
``Random search for hyper-parameter optimization,''
\textit{Journal of Machine Learning Research}, vol. 13, pp. 281--305, 2012.

\bibitem{lorenzo2017psonn}
P. R. Lorenzo, J. Nalepa, M. Kawulok, L. S. Ramos, and J. R. Pastor,
``Particle swarm optimization for hyper-parameter selection in deep neural
networks,''
in \textit{Proc. Genetic and Evolutionary Computation Conference},
pp. 481--488, 2017.

\bibitem{krizhevsky2009cifar10}
A. Krizhevsky,
``Learning multiple layers of features from tiny images,''
Technical Report, University of Toronto, 2009.

\bibitem{kingma2014adam}
D. P. Kingma and J. Ba,
``Adam: A method for stochastic optimization,''
in \textit{Proc. International Conference on Learning Representations},
2015.

\bibitem{natarajan2013noisy}
N. Natarajan, I. S. Dhillon, P. Ravikumar, and A. Tewari,
``Learning with noisy labels,''
in \textit{Advances in Neural Information Processing Systems},
vol. 26, pp. 1196--1204, 2013.

\end{thebibliography}

\end{document}
